\TieudeT{PHONG TỎA VẬT LÍ 12 - VẬT LÍ NHIỆT}
\subsubsection*{PHẦN I. Thí sinh trả lời câu hỏi từ câu 1 đến câu 18. Mỗi câu hỏi thí sinh chỉ chọn một phương án}
\setcounter{ex}{0}
\begin{ex}%chuan
	Phát biểu nào sau đây là \textbf{sai} khi nói về cấu tạo chất?	
	\choice
	{Các phân tử luôn luôn chuyển động không ngừng}
	{Các phân tử chuyển động càng nhanh thì nhiệt độ của vật càng cao}
	{\True Các phân tử luôn luôn đứng yên và chỉ chuyển động khi nhiệt độ của vật cao}
	{Các chất được cấu tạo từ các hạt riêng biệt gọi là phân tử}
	\loigiai{
		\begin{itemize}
			\item Các phân tử chuyển động không ngừng.
			Chuyển động của các phân tử được gọi là chuyển động nhiệt.
		\end{itemize}
	}
\end{ex}


\begin{ex}%chuan
	Lực liên kết giữa các phân tử	
	\choice
	{là lực hút}
	{là lực đẩy}
	{tuỳ thuộc vào thể của nó, ở thể rắn là lực hút còn ở thể khí là lực đẩy}
	{\True gồm cả lực hút và lực đẩy}
	\loigiai{
		\begin{itemize}
			\item Giữa các phân tử có lực hút và lực đẩy, gọi chung là lực liên kết phân tử.
		\end{itemize}
	}
	\end{ex}

\begin{ex}%chuan
	Phát biểu nào sau đây nói về chuyển động của phân tử là \textbf{không} đúng?
	\choice
	{\True Chuyển động của phân tử hoàn toàn là do lực tương tác phân tử gây ra}
	{Các phân tử chuyển động không ngừng}
	{Các phân tử chuyển động nhiệt càng nhanh thì nhiệt độ của vật càng cao}
	{Các phân tử khí chuyển hỗn loạn không ngừng về mọi hướng}
	\loigiai{
		\begin{itemize}
			\item Giữa các phân tử có lực hút và lực đẩy, gọi chung là lực liên kết phân tử.
			\item Lực tương tác phân tử có thể gây thay đổi chuyển động chứ không gây ra chuyển động.
		\end{itemize}
	}

\end{ex}

\begin{ex}%chuan
	Các chất có thể tồn tại ở những thể cơ bản nào?
	\choice
	{Thể rắn}
	{Thể lỏng}
	{Thể khí}
	{\True Thể rắn, thể lỏng, thể khí}
	\loigiai{
		\begin{itemize}
			\item Ba thể cơ bản của các chất: Thể rắn, thể lỏng, thể khí
		\end{itemize}
	}
	\end{ex}

\begin{ex}%chuan
	Vật rắn có hình dạng xác định vì phân tử cấu tạo nên vật rắn
	\choice
	{không chuyển động}
	{đứng sát nhau}
	{chuyển động với vận tốc nhỏ không đáng kể}
	{\True dao động quanh một vị trí cân bằng xác định}
	\loigiai{
		\begin{itemize}
			\item Các phân tử chất rắn được giữ ở các vị trí cân bằng và mỗi phân tử chỉ có thể dao động xung quanh vị trí cân bằng xác định này nên các chất ở thể rắn có thể tích và hình dạng xác định.
		\end{itemize}
	}
	\end{ex}
\begin{ex}%chuan
	Các nguyên tử trong một miếng sắt có tính chất nào sau đây?
	\choice
	{Khi nhiệt độ tăng thì nở ra}
	{Khi nhiệt độ giảm thì co lại}
	{\True Khoảng cách rất gần nhau}
	{Khoảng cách rất xa nhau}
	\loigiai{
		\begin{itemize}
			\item Các nguyên tử trong kim loại nói riêng và chất rắn nói chung thường đứng rất gần nhau.
		\end{itemize}		
	}
	\end{ex}

\begin{ex}%chuan
	Trong tinh thể, các hạt (nguyên tử, phân tử, ion)
	\choice
	{\True dao động nhiệt xung quanh vị trí cân bằng xác định}
	{đứng yên tại những vị trí xác định}
	{chuyển động hỗn độn không ngừng}
	{chuyển động trên quỹ đạo tròn xung quanh một vị trí xác định}
	\loigiai{
		\begin{itemize}
			\item Trong tinh thể, các nguyên tử, phân tử, hoặc ion không đứng yên mà dao động xung quanh các vị trí cân bằng của chúng. 
			\item Vị trí cân bằng này là các vị trí cấu trúc xác định mà các hạt chiếm giữ. Khi nhiệt độ tăng, dao động nhiệt của các hạt tăng lên, nhưng chúng vẫn dao động xung quanh vị trí cân bằng cố định. 
			\item Đây là lý do vì sao tinh thể có hình dạng và cấu trúc ổn định ở điều kiện bình thường.
		\end{itemize}
	}
\end{ex}

\begin{ex}%chuan \nguon{4}
	Khi chất rắn kết tinh được nung nóng. Phát biểu nào sau đây là đúng?
	\choice
	{các phân tử vẫn dao động với biên độ không đổi, khoảng cách trung bình giữa các phân tử không đổi}
	{\True các phân tử dao động với biên độ tăng lên, khoảng cách trung bình giữa các phân tử tăng lên}
	{các phân tử dao động với biên độ không đổi, khoảng cách trung bình giữa các phân tử tăng lên}
	{các phân tử dao động với biên độ tăng lên, khoảng cách trung bình giữa các phân tử không đổi}
	\loigiai{
		Khi nung nóng, các phân tử trong chất rắn dao động mạnh hơn (biên độ tăng), làm cho khoảng cách giữa các phân tử cũng tăng lên.
	}
\end{ex}

\begin{ex}
	"Độ không tuyệt đối" là nhiệt độ ứng với
	\choice
	{\True $0\ K$}
	{$0\ ^\circ C$}
	{$273\ ^\circ C$}
	{$273\ K$}
	
	\loigiai{}
\end{ex}

\begin{ex}
 Chất khí có
	\choice
	{hình dạng xác định}
	{lực tương tác giữa các phân tử rất mạnh}
	{\True lực tương tác giữa các phân tử rất yếu}
	{thể tích riêng xác định}
	\loigiai{}
\end{ex}

\begin{ex}
Hiện tượng nào \textbf{không} liên quan đến sự nóng chảy?
\choice
{Đốt một cây nến}
{\True Thắp một ngọn đèn dầu}
{Đúc một pho tượng đồng}
{Bỏ một cục nước đá vào trong một cốc nước}	
	\loigiai{}
\end{ex}

\begin{ex}
Nhiệt kế rượu \textbf{không thể} đo nhiệt độ nào trong các nhiệt độ sau đây?
\choice
{Nhiệt độ của nước đá}
{Nhiệt độ cơ thể người}
{Nhiệt độ của khí quyển}
{\True Nhiệt độ sôi của nước}	
	\loigiai{}
\end{ex}

\begin{ex}
	Phát biểu nào sau đây là \textbf{sai}? Cùng một khối lượng của một chất nhưng khi ở các thể khác nhau thì sẽ khác nhau về
	\choice
	{thể tích}
	{khối lượng riêng}
	{\True kích thước của các nguyên tử}
	{trật tự của các nguyên tử}
	\loigiai{}
\end{ex}

\begin{ex}
	Phát biểu nào sau đây là đúng?
\choice
{Vận tốc chuyển động của các phân tử có liên quan đến trọng lượng riêng của vật}
{\True Vận tốc chuyển động của các phân tử có liên quan đến nhiệt độ của vật}
{Vận tốc chuyển động của các phân tử có liên quan đến thể tích của vật}
{Vận tốc chuyển động của các phân tử có liên quan đến khối lượng của vật}	
	\loigiai{}
\end{ex}

\begin{ex}
	Phát biểu nào sau đây là đúng khi nói về sự nóng chảy và đông đặc?
	\choice
	{\True Sự nóng chảy là sự chuyển từ thể rắn sang thể lỏng, sự đông đặc là sự chuyển từ thể lỏng sang thể rắn}
	{Sự nóng chảy là sự chuyển từ thể lỏng sang thể lỏng, sự đông đặc là sự chuyển từ thể rắn sang thể rắn}
	{Sự nóng chảy là sự chuyển từ thể rắn sang thể rắn, sự đông đặc là sự chuyển từ thể lỏng sang thể lỏng}
	{Sự nóng chảy là sự chuyển từ thể lỏng sang thể rắn, sự đông đặc là sự chuyển từ thể rắn sang thể lỏng}
	\loigiai{}
\end{ex}

\begin{ex}
Điểm cố định trên và điểm cố định dưới của một nhiệt kế bị hỏng lần lượt là $96^\circ C$ và $-2^\circ C$. Số chỉ của nhiệt kế này ứng với nhiệt độ $50^\circ C$ của nhiệt kế chỉ đúng là
\choice
{$21^\circ C$}
{$12^\circ C$}
{$45^\circ C$}
{\True $47^\circ C$}
	\loigiai{}
\end{ex}

\begin{ex}
	Thả đồng thời $0,2\ kg$ sắt ở $150^\circ C$ và $450\ g$ đồng ở nhiệt độ $25^\circ C$ vào $150\ g$ nước ở nhiệt độ $80^\circ C$. Biết rằng sự hao phí nhiệt vì môi trường là không đáng kể và nhiệt dung riêng của sắt, đồng, nước lần lượt bằng $460\ J/kg.K$, $400\ J/kg.K$ và $4200\ J/kg.K$. Khi cân bằng, nhiệt độ của hệ là
	\choice
	{\True $76,16^\circ C$}
	{$40^\circ C$}
	{$65^\circ C$}
	{$23^\circ C$}
	\loigiai{}
\end{ex}

\begin{ex}
Một cái nồi bằng nhôm chứa nước ở $20^{\circ}C$. Đổ thêm vào nồi $1$ lít nước sôi thì nhiệt độ của nước trong nồi là $45^{\circ}C$. Bỏ qua sự mất mát nhiệt ra môi trường ngoài trong quá trình trao đổi nhiệt bỏ qua trao đổi nhiệt với nồi nhôm. Để nhiệt độ của nước trong nồi là $60^{\circ}C$ thì phải đổ thêm một lượng nước sôi là

\choice
{\True $1,2$ lít}
{$1,5$ lít}
{$1,8$ lít}
{$1,6$ lít}
	\loigiai{}
\end{ex}
\subsubsection*{PHẦN 2. Thí sinh trả lời từ câu 1 đến câu 4. Trong mỗi ý a), b), c), d) ở mỗi câu, thí sinh chọn đúng hoặc sai.}
\setcounter{ex}{0}
\begin{ex}
	Một vận động viên nhảy cầu có khối lượng $m = 55\ kg$ thực hiện động tác nhảy cầu từ độ cao $5\ m$ xuống một bể bơi. Bỏ qua sự trao đổi nhiệt của nước trong bể bơi với môi trường bên ngoài và toàn bộ năng lượng của người thực hiện truyền hết qua nước. Lấy $g = 10\ m/s^2$.
	\choiceTFt
	{Nội năng của nước trong bể bơi thay đổi chủ yếu là do quá trình truyền nhiệt cơ thể vận động viên sang nước trong bể bơi.}
	{Độ biến thiên của nước trong bể bằng độ biến thiên nội năng của cơ thể vận động viên.}
	{Cơ thể vận động viên đã truyền một nhiệt lượng là $2750\ J$ cho bể nước.}
	{\True Độ biến thiên nội năng của nước trong bể bơi là $2750\ J$}
	\loigiai{
		\begin{itemchoice}
			\itemch Do quá trình thực hiện công
			\itemch
			\itemch Do quá trình thực hiện công
			\itemch $\Delta U = mgh = 55.10.5 = 2750\ J$.
		\end{itemchoice}
	}
\end{ex}

\begin{ex}
Một học sinh luộc khoai tây để nấu súp. Học sinh này cho $0,5\ kg$ khoai tây vào nồi nước. Trong quá trình nấu, nhiệt độ của khoai tây tăng từ $20^{\circ}C$ đến $100^{\circ}C$. Biết nhiệt dung riêng của khoai tây là $3,4.10^{3}\ J/kg.K$.
\choiceTFt
{\True Độ biến thiên năng lượng nhiệt của khoai tây bằng $1,36.10^{5}\ J$}
{\True Trong thực tế, năng lượng do bếp cung cấp lớn hơn nhiệt lượng mà khoai tây nhận được do sự tỏa nhiệt ra môi trường xung quanh}
{\True Tăng hiệu suất của nguồn nhiệt thì bạn học sinh có thể giảm thời gian đun khoai tây nóng đến $100^{\circ}C$}
{\True Sau khi đã nấu xong, bạn học sinh cho khoai tây vào máy xay thực phẩm. Máy xay có một động cơ làm quay lưỡi dao để cắt khoai tây. Công suất toàn phần của động cơ là $5,00.10^{2}\ W$. Công suất có ích của động cơ là $3,00.10^{2}\ W$. Hiệu suất của động cơ của máy xay thực phẩm là $60\%$}
	\loigiai{
		\begin{itemchoice}
			\itemch Độ biến thiên năng lượng nhiệt của khoai tây bằng nhiệt lượng mà nó nhận được: $Q = m.c.\Delta t = 0,5.3,4.10^3(100-20) = 1,36.10^5\ J$.
			\itemch 
			\itemch Tăng hiệu suất bằng cách: đậy nắp nồi, sử dụng tấm chắn gió,...
			\itemch $ H = \dfrac{P_{ci}}{P_{tp}} = 60\%$.
		\end{itemchoice}
	}
\end{ex}

\begin{ex}
	Một xe máy có công suất $3,2\ kW$ đang chuyển động với vận tốc $45\ km/h$. Biết hiệu suất của động cơ là $25\%$, năng suất toả nhiệt của xăng là $4,6.10^7\ J/kg$, khối lượng riêng của xăng là $700\ kg/m^3$. 
	\choiceTFt
	{Nhiệt lượng tỏa ra khi đốt hết $1$ lít xăng là $4,6. 10^7\ J$}
	{\True Nếu xe tiêu thụ hết $1$ lít xăng thì động cơ đã thực hiện công là $8,05\ MJ$}
	{\True Trung bình mỗi lít xăng tiêu thụ xe sẽ đi được quãng đường khoảng $31\ km$}
	{Khi đi $45\ km$ thì động cơ đã thực hiện một công là $12.10^7\ J$}
	\loigiai{
		\begin{itemchoice}
			\itemch 
			\begin{itemize}
				\item $m = V.D = 10^{-3}.700 = 0,7\ kg$.
				\item $Q = q.m = 4,6.10^7.0,7 = 3,22.10^7\ J$.
			\end{itemize}
			\itemch $A = H.Q = 0,25.3,22.10^7 = 8,05.10^6\ J = 8,05\ MJ$.
			\itemch $s=v.t = v.\dfrac{A}{P} = 45.\dfrac{8,05.10^6}{3,2.10^3.3600} \approx 31,45\ km$.
			\itemch $A_d = P.t_d = 3,2.10^3.3600 = 11,52.10^6\ J$.
		\end{itemchoice}
	}
\end{ex}

\begin{ex}
	Một học sinh làm thí nghiệm đun nóng để làm $0,02\ kg$ nước đá (thể rắn) ở $0\ ^{\circ}C$ chuyển hoàn toàn thành hơi nước ở $100\ ^{\circ}C$. Cho nhiệt nóng chảy của nước ở $0\ ^{\circ}C$ là $3,34.10^{5}\ J/kg$; nhiệt dung riêng của nước là $4,2\ kJ/kgK$; nhiệt hoá hơi riêng của nước ở $100\ ^{\circ}C$ là $2,26.10^{6}\ J/kg$. Bỏ qua hao phí toả nhiệt ra môi trường.
	\choiceTFt
	{Nhiệt lượng cần thiết để làm nóng chảy hoàn toàn $0,02\ kg$ nước đá tại nhiệt độ nóng chảy là $6860\ J$}
	{Nhiệt lượng cần thiết để đưa $0,020\ kg$ nước từ $0\ ^{\circ}C$ đến $100\ ^{\circ}C$ là $8600\ J$}
	{Nhiệt lượng cần thiết để làm hoá hơi hoàn toàn $0,020\ kg$ nước ở $100\ ^{\circ}C$ là $42500\ J$}
	{\True Nhiệt lượng để làm $0,020\ kg$ nước đá (thể rắn) ở $0\ ^{\circ}C$ chuyển hoàn toàn thành hơi nước ở $100\ ^{\circ}C$ là $60280\ J$}
	\loigiai{
		\begin{itemchoice}
			\itemch $Q_1 = \lambda .m = 3,34.10^5.0,02 = 6680\ J$.
			\itemch $Q_2 = m.c.\Delta t = 0,02.4200.100 = 8400\ J$.
			\itemch $Q_3 = L.m = 2,26.10^6.0,02 = 45200\ J$.
			\itemch $Q = Q_1 +Q_2 +Q_3 = 60280\ J$.
		\end{itemchoice}
	}
\end{ex}
\subsubsection*{PHẦN 3. Thí sinh trả lời từ câu 1 đến câu 6.}
\setcounter{ex}{0}
\begin{ex}% Câu 1
	Khoảng cách cơ bản tức là số khoảng chia giữa hai điểm cố định trên và cố định dưới của một thang đo nào đó. Hai thang đo X và Y có khoảng cách cơ bản lần lượt là $50$ và $150$. Điểm đông đặc của nước trên cả hai thang đo này đều là $0^\circ$. Nếu nhiệt độ trên thang đo X là $15^\circ$ thì nhiệt độ tương ứng trên thang đo Y bằng bao nhiêu (làm tròn kết quả đến chữ số hàng đơn vị)?
	\shortans[oly]{$45$}
	\loigiai{
	\begin{itemize}
		\item $\dfrac{\Delta t}{\Delta \ell} = const \Rightarrow \dfrac{15-0}{50-0} = \dfrac{y-0}{150-0} \Rightarrow y = 456^\circ Y$.
	\end{itemize}
	}
\end{ex}

\begin{ex}% Câu 2
	Người ta thả một chai sữa của trẻ em vào phích đựng nước ở nhiệt độ $t = 60^\circ C$. Sau khi đạt cân bằng nhiệt, chai sữa nóng tới nhiệt độ $t_1 = 40^\circ C$, người ta lấy chai sữa này ra và tiếp tục thả vào phích một chai sữa khác giống như chai sữa trên. Biết rằng trước khi thả vào phích, các chai sữa đều có nhiệt độ $t_0 = 20^\circ C$. Xem tỏa nhiệt ra môi trường là không đáng kể. Khi cân bằng nhiệt thì nhiệt độ của chai sữa bằng bao nhiêu?
	\shortans[oly]{$30$}
	\loigiai{
	\begin{itemize}
		\item Khi thả chai sữa thứ nhất: $m_nc_c(t-t_1) = m_sc_s(t_1-t_0) \Rightarrow m_nc_n.20=m_sc_s.20 \Rightarrow m_nc_n = m_sc_s$.
		\item Khi thả chai sữa thứ hai: $m_nc_n(40-t_{cb}) = m_sc_s(t_{cb} - 20) \Rightarrow t_{cb} = 30^\circ C$.
	\end{itemize}
	}
\end{ex}

\begin{ex}% Câu 3
	Có hai bình cách nhiệt: Bình (1) chứa khối lượng $m_1 = 3\ kg$ nước ở nhiệt độ $30^\circ C$, bình (2) chứa khối lượng $m_2 = 5\ kg$ nước ở $70^\circ C$. Người ta rút một lượng nước có khối lượng $m$ từ bình (1) sang bình (2). Sau khi cân bằng nhiệt, người ta lại rút từ bình (2) sang bình (1) một lượng nước có khối lượng cũng bằng $m$. Nhiệt độ cân bằng ở bình (1) là $31,95^\circ C$. Bỏ qua sự trao đổi nhiệt khi rút nước từ bình nọ sang bình kia và giữa nước với bình. Nhiệt độ cân bằng của nước ở bình (2) sau khi rút nước từ bình (1) bằng bao nhiêu $^\circ C$ (làm tròn kết quả đến chữ số hàng đơn vị)?
	
	\shortans[oly]{$68,8$}
	\loigiai{
		\begin{itemize}
			\item Cân bằng nhiệt ở bình 2: $$mc(t-t_1) = m_2c(t_2-t)\Rightarrow m = \dfrac{m_2(t_2-t)}{(t-t_1)} = \dfrac{5(70-t)}{t-30}$$.
			\item Cân bằng nhiệt ở bình 1: $$(3-m)c(t'-t_1) = mc(t-t') \Rightarrow \left( 3- \dfrac{5(70-t)}{t-30}\right) (31,95 - 30) = \dfrac{5(70-t)}{t-30}(t-31,95) \Rightarrow t \approx 68,8^\circ C$$.
		\end{itemize}
	}
\end{ex}

\begin{ex}% Câu 4
Một bác thợ rèn nhúng một con dao bằng thép có khối lượng $1,1\ kg$ ở nhiệt độ $850^\circ C$ vào trong bể nước lạnh để làm tăng độ cứng của lưỡi dao. Nước trong bể có thể tích là $50$ lít và có nhiệt độ bằng với nhiệt độ ngoài trời là $27^\circ C$. Phần hao phí do sự truyền nhiệt cho thành bể và môi trường ngoài chiếm $10\%$ nhiệt lượng do dao tỏa ra. Biết nhiệt dung riêng của thép là $460\ J/kgK$, của nước là $4200\ J/kgK$; khối lượng riêng của nước là $1$ kg/lít. Nhiệt độ của nước khi có sự cân bằng nhiệt là bao nhiêu độ C (làm tròn kết quả đến chữ số hàng phần mười)?
	\shortans[oly]{$28,8$}
	\loigiai{
		\begin{itemize}
			\item $Q_n + Q_{hp} = Q_{\text{thép}} = Q_n + 0,1Q_{\text{thép}} = Q_{\text{thép}} \Leftrightarrow Q_n = 0,9Q_{\text{thép}} \Leftrightarrow m_n.c_n.\Delta t_n = 0,9.m_{\text{thép}}.\Delta t_{\text{thép}} \Rightarrow 50.4200.(t-27) = 0,9.1,1.460.(850-t) \Rightarrow t \approx 28,8$.
		\end{itemize}
	}
\end{ex}

\begin{ex}% Câu 5
	Bỏ $100\ g$ nước đá ở $t_1 = 0^\circ C$ vào $300\ g$ nước ở $t_2 = 20^\circ C$. Cho nhiệt nóng chảy riêng của nước đá là $\lambda = 3,4.10^5\ J/kg$ và nhiệt dung riêng của nước là $c = 4200\ J/kg.K$. Khối lượng nước đá còn lại bằng bao nhiêu gam (làm tròn kết quả đến chữ số hàng đơn vị)?
	\shortans[oly]{$26$}
	\loigiai{
		\begin{itemize}
			\item Nhiệt lượng nước tỏa ra: $Q_2 = m_2c\Delta t_2 = 0,3.4200.20=25200\ J$.
			\item Khối lượng đá tan là: $m_{tan} = \dfrac{Q_2}{\lambda} = \dfrac{25200}{3,4.10^5} = \dfrac{63}{850}\ kg$.
			\item Khối lượng nước đá còn lại: $m - m_{tan} = 100 - \dfrac{63}{850}.1000 \approx 26\ g$.
			
		\end{itemize}
	}
\end{ex}

\begin{ex}% Câu 6
	Để biến $500\ g$ nước ở $30^\circ C$ thành nước đá, người ta bỏ vào nước trên một khối nước đá ở $-10^\circ C$. Biết nhiệt dung riêng của nước là $c_1 = 4200\ J/kg.K$ và nước đá là $c_2 = 2000\ J/kg.K$; nhiệt nóng chảy riêng của nước đá $\lambda = 3,4.10^5 \ J/kg$. Lượng nước đá tối thiểu cần dùng bằng bao nhiêu kg (làm tròn kết quả đến chữ số hàng đơn vị)?
	
	\shortans[oly]{$12$}
	\loigiai{
	\begin{itemize}
		\item Nhiệt lượng của $500\ g$ nước ở $30^\circ C$ tỏa ra để thành nước đá là:
		$$Q = m_1c_1\Delta t +\lambda m_1 = 0,5.4200.30 + 3,4.10^5.0,5 = 233000\ J$$.
		\item Nhiệt lượng này đúng bằng nhiệt lượng thu vào của nước đá tăng từ $-10^\circ C$ đến $0^\circ C$:
		$$Q = m_2c_2\Delta t_2 \Rightarrow m_2 = \dfrac{Q}{c_2\Delta t_2} = \dfrac{233000}{2000.10} = 11,65 \approx 12\ kg$$.
	\end{itemize}
	}
\end{ex}
